%/*% pdflatex -file-line-error -interaction nonstopmode %

\documentclass{article}
\usepackage{amsmath}
\usepackage{amsfonts}
\usepackage{hyperref}
\usepackage{exercises}
\usepackage{amsthm}
\newcommand*{\oldneg}{\mathord{\sim}}

\theoremstyle{definition}
\newtheorem{definition}{Definition}

\setlength{\parindent}{0pt}
\setlength{\parskip}{10pt}

\title{Chapter 11}
\begin{document}
\maketitle
We are attempting to prove that every decimal number has a unique binary and dyadic representation.
\section*{Uniqueness}
Dyadic Notation Case:\newline
\begin{proof}
We assume that there exists at least two dyadic representations for some decimal number.
\newline Result 1: Consider any number $2\times 2^n$. We are to show that this number is bigger than any sum of any $2^m, 2\times 2^m$ for $m<n$ where every
term with every power $i$ is counted only once. In the base case, we have $2\times 2^2=8$, while the greatest possible sum of lower $m$'s is
$2\times 2^1+2\times 2^0=6$, thus completing this case. Make an induction assumption. Well, consider $2\times 2^{n+1}$. 
We can rewrite this as $2^{n+2}$. Then,consider the greatest possible sum of lower powers of $m$, that is:\newline
$2\times 2^n+2\times 2^{n-1}+...+2\times 2^0$. By the assumption, this simplifies to:\newline
$\rightarrow 2\times 2^n+m$ where $m<2\times 2^n$. This shows that the sum is:
$2^{n+1}+m$, that is, some number less than $2^{n+1}+2^{n+1}$, meaning the sum is less than $2^{n+2}$, thus completing this part of the proof.\newline
\newline Final Result: There are then two cases for this non-unique dyadic representation. Either the two representations are the same length, or they are
not. \newline Consider the first case, where they are of the same length. Well, because they are different representations, find the first spot where the
representations differ. Hence, at this point, one representation has a $1$, while the other has $2$. The one with the $2$ at this points currently
represents a number greater than the one that the other representation represents. For these to be equal, the next lowest digit must be a $2$ for the smaller
representation. However, the greater one stays greater, as it must at least have a $1$ for this digit. Hence, this process must continue if they are indeed 
equal decimals. Well, at no point can they be equal, as, whenever we add a $2$ to the lesser representation, a $1$ is added to the greater one, thus 
meaning it is still larger, and hence, the representations do not represent the same decimal number.\newline Now Consider the second case where the
representations are different lengths. If the longer representation has a $2$ in the $2^n$ place for the largest $n$, the decimal numbers that the
representations represent are not equal by Result 2. Otherwise, the longer representation has a $1$ in the $2^n$ place. Well, for these two representations
to be equal, the $2^m$ digit for the shorter representation for the greatest $m$ must be $2$, while the digit in the $m$ place in the longer representation
must be $1$. This process must go on, and, as we saw in the first case, it never ends with the representations being equal (because the greater
representation is of a longer length), thus proving they represent different decimal numbers.
\end{proof} 
Binary Notation Case:\newline
\begin{proof}
We temporarily assume that some binary notation is non-unique for some number.\newline
\newline Result 1: Consider any number $2^n$. We are to show that no sum of any number $2^m$ for $m<n$ where every $2^i$ is counted at most once, 
can be equal to or greater than $2^n$. In the base case of $1$, this is clear, as the only possible sum of lower $2^i$'s would be $2^0$, that is, $1$, 
which does not equal $2^1=2$. Now, assume that the statement holds true for some $n$. Now, consider $2^{n+1}$. Now, to show that all possible sums of the 
manner that we are considering are less than $2^n+1$, we need only to show that the greatest sum, namely the one in which all $n-1$ lower powers of $2$ are 
counted once, is less than $2^{n+1}$. Well, consider this sum. We consider the partial sum $2^0+...+2^{n-1}$. By the assumption, this is less 
than $2^n$. Hence, $2^n+2^{n-1}+...+2^0=2^n+m$ where $m<2^n$. Thus, $2^n+2^{n-1}+...+2^0<2^n+1$, thus completing this inductive proof.\newline
\newline Final Result: There are then two cases: either the lengths of the binary numbers are the same, or it is not.
\newline \newline Consider the first case. Starting at $2^n$, subtract
$2^n$, and then, if the lengths of the resulting numbers are the same, subtract $2^{n-1}$, and so on. If we never reach a point where the numbers differ, the
binary notations are clearly the same. However, if we ever reach some point where the digit for one number is $2^i$ is $1$, while the other is $0$, no binary
number up to length $i-1$ can equal $2^i$ by Result 1 (and it must equal $2^i$ as we assumed that we had subtracted equivalent numbers at every earlier
stage, thus preserving equality). Hence, these representations do not represent the same decimal number.\newline \newline 
Consider the second case. Well, this implies that the digit in the $2^m$ place is $1$ for $m$ the length of the longer binary number. Well, by Result 2, 
no possible sum of any length $n<m$ could equal $2^m$, thus proving that these two numbers do not represent the same decimal number, thus completing 
the proof. 
\end{proof}
\section*{Completeness}
Binary Notation Case:\newline
\begin{proof}
Result 1: \newline We first prove an important Lemma, that is, every integer is either a power of $2$, or equals a sum of distinct powers of $2$.
In the base case, $1$ is a power of $2$, thus completing the base case. Now, make a complete induction assumption.There are two cases: either the integer
is odd, or it is even. Let us first consider the case where the integer is even. So, let us consider the $n+1$th even
number. Well, by the assumption, the $n$th even number is either a power of $2$ or equals a sum of powers of $2$. Now, let us consider $m$, the $n$th even
number. If it is a power of $2$, it must be $\ge 2$, as we assumed that it was even. Well, if $m=2$, this means that $n$, the $n+1$th even number, is
$2+2=4$, and is therefore a power of $2$. Otherwise, $m$ is some greater power of $2$, thus meaning that $2+m$, a sum of two distinct powers of $2$, equals
$n$, thus completing this case.\newline In the other case, we have $m$ equalling a sum of powers of $2$. Let us consider this sum. If $2$ is not a term of
the sum, then $2+m=n$, that is, some sum of distinct powers of $2$ $=n$. Otherwise, it contains $2$. So, take $2$ out of the sum, and pair it with the $2$
in the equation $2+m=n$. Then, the leftover sum for $m$ either has $4$ or it does not. If it does not, we have $4+k=n$ for $k$ the sum for $m$ without $2$.
Otherwise, we pair the $4$ with the leftover $4$ in the equation $4+k=n$, and we repeat this process. Eventually, this process must end. When the process
ends, we would be left with $2^l+s=n$ for $l$ the number of times the process repeated, and $s$ the leftover partial sum for $m$, thus meaning $n$ is equal
to a sum of distinct powers of $2$, thus completing the proof. In the case where the integer is odd, we simply take the sum of powers of $2$ for the
odd integer minus one, and then add to this sum the term $2^0=1$.
\newline Final Result:\newline Consider any decimal number $n$. It is either $0$, or it is not. If it equals $0$, the binary number $0$ represents this. 
Otherwise, we consider the sum of distinct powers of $2$ that $n$ has (by Result 1). To make a binary representation for this decimal, we simply
pick a digit of $1$ for all those powers of $2$ listed in its corresponding sum of powers of $2$, and pick a $0$ for any other powers of $2$, thus 
effectively representing this number in binary notation.
\end{proof}
Dyadic Notation Case:\newline
\begin{proof}
Consider any decimal number $n$. Now, let us consider its binary representation shown in the above proof. Here is how we can translate this binary 
notation into dyadic notation. The first time we encounter a $0$ in $k$th digit in the binary notation, we shift over the representation by firsting changing
the $k$th digit to a $2$ and making the $k+1$th digit a $0$. Then, we go to the $k+1$th digit, and repeat this process until we reach the last digit. Then,
 whenever we encounter another $0$ in the rest of the notation, if the $k+1$th digit is a $1$, we change the $k$th digit to a $2$ and the $k+1$th digit to a 
$0$. If the $k+1$th digit is a $2$, we change the $k+1$th digit to a $1$, and the $k$th digit to a $2$. If the $k+1$th digit was a $1$, we repeat the process
 until reach the last digit. It should be clear that the process described conserves the initial decimal number that the representation represented. Also, 
because we repeat this process until there are no more $0$'s, it is clear this process removes all $0$'s from the notation. Thus, we have effectively gone 
from binary notation to a notation with only $1$'s and $2$'s which represents the same decimal number, and hence, we have therefore represented the decimal
 number in dyadic notation. Of course, if the binary notation is $only$ $0$'s, then the decimal number it represents is $0$, and thus we need not worry about
 this case, as we are only trying to show all $integers$ are representable in dyadic notation.
\end{proof}
\section*{A Further Result}
Proof the every dyadic or binary string of $n$ $1$'s is equal to $2^n-1$.
\begin{proof}
In the base case, one $1=2^0=1$, that is, $2^1-1$. Now, make an induction assumption. So, a string of $n+1$ $1$'s is equal to $1\times 2^n + 2^n-1$ that is,
 $2^{n+1}-1$.
\end{proof} 
\end{document}
